% Drawing 152/01-1 -- folder 152, batch 01, archive page 1 (PDF page 2).
% Transcription: en haut à gauche du feuillet : un cercle portant quatre sommets marqués, de chacun desquels part vers l'extérieur un faisceau de trois traits ou plus (celui de droite très chargé) ; autour, quatre flèches courbes indiquant un sens de rotation.
% Kind: circle with four marked vertices, bundles of rays and rotation arrows. Method 2 (schematic redraw from the facsimile). Status: ok.
% Notes: The circle, four marked vertices on it (top, left, right, bottom), and from each a bundle of strokes going outwards, all drawn from the vertex: three at the top, two at the left, six at the right (three heavy ones close together to the upper right, three thinner ones to the right and lower right), three at the bottom. Round the figure four curved arrows: above (pointing left), at the left (pointing down), at the right (pointing up), below (pointing right). Overlaid on the 400 dpi colour facsimile on 2026-11-12 (Opus 5.5) and corrected: the heavy bundle at the right is two long heavy strokes and the blot between them (drawn as three heavy strokes), and the thinner rays are three (upper right, right, lower right); the short seventh ray of the earlier redraw, a doubling of the lower right one, removed; the faint short stroke inside the circle at the left added (pencil). The four arrows all turn counter-clockwise round the circle (top to the left, left downwards, bottom to the right, right upwards), as drawn. He went round the circle twice at the lower left and at the bottom; it is drawn once. Redrawn 2026-11-12 (issue 128): to check again.
% Coordinates: 1 unit = 1 mm = 1/100 of the crop width (bbox in page fractions: [0.12, 0.07, 0.3, 0.195]).
% Redrawn by Claude (claude-fable-5-2), issue #128, 2026-11-12; not checked by a human; not compiled (no TeX in the session).
% Corrected by Claude (claude-opus-5-5) 2026-11-12 (issue #128): overlaid on the 400 dpi colour facsimile and corrected from it. Not compiled. Status: ok.
% Rights: reproduces part of the work -- behind RENDER_TRANSCRIPTS, never published by itself.
\documentclass[tikz,border=3pt]{standalone}
\usepackage{amsmath,amssymb}
\usetikzlibrary{patterns,arrows.meta}
\tikzset{pencil/.style={black!45}, struck/.style={black!35}, mpt/.style={circle,draw,inner sep=1.3pt}, pt/.style={circle,fill,inner sep=0.9pt}, opt/.style={circle,draw,fill=white,inner sep=1.1pt}}
\begin{document}
\begin{tikzpicture}[x=1mm,y=1mm,line cap=round,line join=round,font=\scriptsize]
\draw (50.2,-48) circle[radius=16.6];
\node[pt] at (55.3,-30.4) {};
\draw[pencil] plot[smooth, tension=0.55] coordinates {(39.4,-36.6) (40.3,-39.2) (40.4,-42.6)};
\node[pt] at (33.9,-46.8) {};
\node[pt] at (66.3,-50.6) {};
\node[pt] at (45.1,-64.1) {};
\draw (55.3,-30.4) -- (50.2,-17.4);
\draw (55.3,-30.4) -- (60.4,-8.4);
\draw (55.3,-30.4) -- (72.5,-11.5);
\draw (33.9,-46.8) -- (16.2,-40.9);
\draw (33.9,-46.8) -- (17,-55.2);
\draw[very thick] (66.3,-50.6) -- (80.4,-30.2);
\draw[very thick] (66.3,-50.6) -- (78.3,-35.3);
\draw[very thick] (66.3,-50.6) -- (77.6,-38.6);
\draw (66.3,-50.6) -- (83.4,-40.4);
\draw (66.3,-50.6) -- (82,-51.6);
\draw (66.3,-50.6) -- (82.2,-62.1);
\draw (45.1,-64.1) -- (35.7,-81.8);
\draw (45.1,-64.1) -- (50.2,-82.5);
\draw (45.1,-64.1) -- (57.1,-75.6);
\draw[-{Stealth[length=2mm]}] plot[smooth, tension=0.55] coordinates {(66.8,-5.4) (60.4,-4.1) (52.8,-5.9)};
\draw[-{Stealth[length=2mm]}] plot[smooth, tension=0.55] coordinates {(9.3,-37.3) (7.3,-44.2) (8.3,-50.9)};
\draw[-{Stealth[length=2mm]}] plot[smooth, tension=0.55] coordinates {(89.8,-57) (92.9,-46.8) (91.9,-35.8)};
\draw[-{Stealth[length=2mm]}] plot[smooth, tension=0.55] coordinates {(40,-88.7) (50.2,-88.4) (58.4,-85.1)};
\end{tikzpicture}
\end{document}
